\documentclass[11pt]{article}

\usepackage{hw_style}

% Set due date for this homework (updated to 2025)
\setduedate{09}{23}

% Setup homework number and header
\hwnumber{0}

\title{
    \vspace{-1.5cm}
    {\Huge\color{headercolor}\textbf{数学物理方法}}\\
    \vspace{0.3cm}
    {\color{accentcolor}\rule{8cm}{1pt}}\\
    \vspace{0.3cm}
    {\Large\color{sectioncolor}温故知新}
}
\author{
    \begin{tabular}{rl}
        \textbf{\color{sectioncolor}任课教师:} & 罗凯\\[0.2cm]
    \end{tabular}
}


\date{2025-9-16}



\begin{document}


\maketitle


\begin{hwproblem}{1}
求解下列问题：
\begin{enumerate}
  \item 方程$x^3 -3x^2 + 2x = 0 $的解；
  \item 两直线（1）$3 x+4 y-2=0$，（2）$2 x+2 y+2 = 0$的交点坐标;
  \item 写出椭圆方程的标准形式；
  \item 写出向量$\vec{a},\vec{b}$夹角的公式；
  \item 四个数$3,3,7,7$算24点。
\end{enumerate}
\end{hwproblem}

\vspace{0.5cm}

\begin{hwproblem}{2}
  写出下列公式：
  \begin{enumerate}
    \item 首项为$a_1$,公差为$d$的等差数列求和公式；
    \item 首项为$a_1$,公比为$q$的等比数列求和公式；
    \item 和差化积公式 $\sin \alpha + \sin \beta$；
    \item 和差化积公式 $\cos \alpha - \cos \beta$；
    \item 积化和差公式 $\sin\alpha \cos\beta$；
    \item 积化和差公式 $\sin\alpha \sin \beta$。
  \end{enumerate}
\end{hwproblem}

\vspace{0.5cm}

\begin{hwproblem}{3}
  求下列极限:
  \begin{enumerate}
    \item $\lim _{n \rightarrow \infty} \frac{n^2+4 n+3}{2 n^2}$；
    \item $\lim _{x \rightarrow+\infty}\left(\sqrt{x^2+x+1}-\sqrt{x^2-x+1}\right)$；
    \item $\lim _{x \rightarrow 0} \frac{\sin 2 x }{x} $；
    \item $\lim _{n \rightarrow \infty}\left(1+\frac{1}{3n}\right)^n$。
  \end{enumerate}  
\end{hwproblem}

\vspace{0.5cm}

\begin{hwproblem}{4}
  试解下列问题：
  \begin{enumerate}
    \item 求不定积分 $\int \frac{1}{1+x^2} dx$,
    \item 求解函数的导数 $g(x)=\int_0^x \frac{1}{1+t^4} d t$,
    \item 求解函数的导数 $g(x)=\int_0^{\sin x} \sqrt{t+\sqrt{t}} d t$.
  \end{enumerate}
\end{hwproblem}

\vspace{0.5cm}

\begin{hwproblem}{5}
  试解下列问题：
  \begin{enumerate}
    \item 求 $\nabla \cdot \vec{r}$,
    \item 求 $\nabla \times \vec{r}$,
    \item 写出 $\nabla^2(\varphi \psi)$的完整表达式.
  \end{enumerate}
\end{hwproblem}

\begin{hwproblem}{6}
  已知矩阵
  \[
    A = \begin{pmatrix}
      1 & 2 \\
      3 & 4
    \end{pmatrix}
  \]
  求矩阵$A$的行列式和逆矩阵。
\end{hwproblem}
\vspace{0.5cm}

% \hwrequirements

\end{document}